\section{Minimal inputs of the HPA--\texorpdfstring{$\Omega$}{Omega} framework}
\label{sec:minimal_inputs}

This section records the structural modules that will be repeatedly invoked. Proofs and additional discussion are provided in the companion sources \cite{Ma2025HPA,Ma2025Omega}.

\subsection{Multiplicative ontology and polar embedding}
HPA takes the multiplicative monoid of positive integers $M=(\NNpos,\cdot)$ as a primitive structure and defines a polar embedding into $\CC^*$ of the form
\begin{equation}
\mathcal{Z}(n)=\rho(n)\,\e^{\iu\theta_\times(n)}\in \CC^*,
\end{equation}
where $\rho(n)>0$ and $\theta_\times$ is additive under multiplication (mod $2\pi$), so that $\mathcal{Z}$ is a multiplicative homomorphism.

\subsection{A no-go principle: addition and multiplication cannot both be preserved}
HPA uses a rigidity fact: except for the trivial real embedding, there is no embedding of $\NNpos$ into $\CC^*$ that preserves both multiplication and linear addition. One convenient formulation is the following.

\begin{lemma}[Rigidity of additive--multiplicative embeddings on $\NNpos$]
\label{lem:no_go_add_mult}
Let $f:\NNpos\to\CC^*$ satisfy
\begin{equation}
f(m+n)=f(m)+f(n),\qquad f(mn)=f(m)f(n),
\end{equation}
for all $m,n\in\NNpos$. Then $f(n)=n$ for all $n\in\NNpos$.
\end{lemma}
\begin{proof}
By additivity, $f(n)=nf(1)$ for all $n$. By multiplicativity, $f(1)=f(1\cdot 1)=f(1)^2$. Since $f(1)\in\CC^*$, this implies $f(1)=1$, hence $f(n)=n$.
\end{proof}

Therefore any nontrivial polar embedding with a genuine phase component cannot be additive, and ``addition'' can only enter through a readout projection/encoding. The resulting mismatch defines a structural gap that enters the observable layer.

\subsection{The scan operator \texorpdfstring{$\Theta$}{Theta}: time as iteration count}
To decouple ontology (phase/multiplicative structure) from readout (additive projection), HPA introduces a genuine unitary scan operator $\Theta$ and defines time as the iteration count of $\Theta$. The scan shift and the phase-multiplication operator form a Weyl pair \cite{Hall2013,Folland1989}, producing a noncommutative structure and uncertainty-type tradeoffs.

Concretely, in standard quantum kinematics one models a Weyl pair by unitary operators $U,V$ satisfying
\begin{equation}
UV=\e^{\iu\omega}VU,
\end{equation}
for some $\omega\in\RR$. We encode the HPA scan--phase interface in the same form.

\begin{axiom}[Scan--phase Weyl relation]\label{ax:scan_phase_weyl}
There exist unitary operators $\Theta$ (scan shift) and $\Phi$ (phase multiplication) acting on the readout Hilbert space such that
\begin{equation}
\Theta\Phi=\e^{\iu\omega}\Phi\Theta
\end{equation}
for some $\omega\in\RR$ fixed by the scan slope/branch.
\end{axiom}

\subsection{Golden branch and minimal binary readout}
In the golden branch, the scan slope is selected such that the induced binary mechanical word becomes the Fibonacci word; the Ostrowski numeration degenerates to the Zeckendorf representation \cite{MorseHedlund1940,Lothaire2002,Zeckendorf1972}.

The motivation for the golden slope is the ``most irrational'' property of $\varphi$: its continued fraction has the minimal partial quotients, which makes rational approximations maximally difficult \cite{Khinchin1964}. For Kronecker sequences $\{n\alpha\}$, discrepancy and gap statistics are controlled by continued-fraction data \cite{KuipersNiederreiter1974}. In particular, the three-distance theorem implies that, for any irrational $\alpha$, the gaps between successive elements of $\{n\alpha\}$ take at most three values; the detailed ``gap/step'' structure is classical \cite{Slater1967}. These external results justify treating the golden branch as a canonical low-complexity choice for binary readout localization.
